
Four research questions are addressed:

\begin{itemize}
\item RQ1: Does k=5 self-contained tier stratification produce n >= 20 problems per tier per model-benchmark pair?
\item RQ2: Are difficulty tier assignments consistent across architectures (Jaccard > 0.30)?
\item RQ3: Does P(True) extraction produce non-degenerate confidence distributions (std(c) > 0.05)?
\item RQ4: Does DELTA\_ECE differ by architecture, and does global temperature scaling correct the pattern?
\end{itemize}

\textbf{Hardware.} Solution generation on a single NVIDIA H100 GPU (approximately 4h28m for k=5 x 542 problems per model). P(True) extraction approximately 4 minutes for 5,730 pairs on H100. ECE computation is CPU-only.

\textbf{Solution generation.} HuggingFace transformers, temperature=0.8, top\_p=0.95, max\_new\_tokens=512, k=5 independent solutions per problem.

\textbf{ECE computation.} M=15 equal-width bins; bootstrap n=1000 samples (seed=42); 95\% CI as 2.5th-97.5th bootstrap percentiles.

\textbf{Gate criteria (pre-registered).} DELTA\_ECE >= 0.03 AND CI lower bound > 0 in at least 2/3 model families.

\textbf{Null baseline.} Monte Carlo Bernoulli null model (n\_sim=100,000): confidence drawn from the model's empirical distribution with correctness assigned independently. Tests whether observed DELTA\_ECE exceeds what would be expected if confidence and correctness were uncorrelated.
